% Options for packages loaded elsewhere
% Options for packages loaded elsewhere
\PassOptionsToPackage{unicode}{hyperref}
\PassOptionsToPackage{hyphens}{url}
\PassOptionsToPackage{dvipsnames,svgnames,x11names}{xcolor}
%
\documentclass[
  letterpaper,
  DIV=11,
  numbers=noendperiod]{scrartcl}
\usepackage{xcolor}
\usepackage{amsmath,amssymb}
\setcounter{secnumdepth}{-\maxdimen} % remove section numbering
\usepackage{iftex}
\ifPDFTeX
  \usepackage[T1]{fontenc}
  \usepackage[utf8]{inputenc}
  \usepackage{textcomp} % provide euro and other symbols
\else % if luatex or xetex
  \usepackage{unicode-math} % this also loads fontspec
  \defaultfontfeatures{Scale=MatchLowercase}
  \defaultfontfeatures[\rmfamily]{Ligatures=TeX,Scale=1}
\fi
\usepackage{lmodern}
\ifPDFTeX\else
  % xetex/luatex font selection
\fi
% Use upquote if available, for straight quotes in verbatim environments
\IfFileExists{upquote.sty}{\usepackage{upquote}}{}
\IfFileExists{microtype.sty}{% use microtype if available
  \usepackage[]{microtype}
  \UseMicrotypeSet[protrusion]{basicmath} % disable protrusion for tt fonts
}{}
\makeatletter
\@ifundefined{KOMAClassName}{% if non-KOMA class
  \IfFileExists{parskip.sty}{%
    \usepackage{parskip}
  }{% else
    \setlength{\parindent}{0pt}
    \setlength{\parskip}{6pt plus 2pt minus 1pt}}
}{% if KOMA class
  \KOMAoptions{parskip=half}}
\makeatother
% Make \paragraph and \subparagraph free-standing
\makeatletter
\ifx\paragraph\undefined\else
  \let\oldparagraph\paragraph
  \renewcommand{\paragraph}{
    \@ifstar
      \xxxParagraphStar
      \xxxParagraphNoStar
  }
  \newcommand{\xxxParagraphStar}[1]{\oldparagraph*{#1}\mbox{}}
  \newcommand{\xxxParagraphNoStar}[1]{\oldparagraph{#1}\mbox{}}
\fi
\ifx\subparagraph\undefined\else
  \let\oldsubparagraph\subparagraph
  \renewcommand{\subparagraph}{
    \@ifstar
      \xxxSubParagraphStar
      \xxxSubParagraphNoStar
  }
  \newcommand{\xxxSubParagraphStar}[1]{\oldsubparagraph*{#1}\mbox{}}
  \newcommand{\xxxSubParagraphNoStar}[1]{\oldsubparagraph{#1}\mbox{}}
\fi
\makeatother


\usepackage{longtable,booktabs,array}
\usepackage{calc} % for calculating minipage widths
% Correct order of tables after \paragraph or \subparagraph
\usepackage{etoolbox}
\makeatletter
\patchcmd\longtable{\par}{\if@noskipsec\mbox{}\fi\par}{}{}
\makeatother
% Allow footnotes in longtable head/foot
\IfFileExists{footnotehyper.sty}{\usepackage{footnotehyper}}{\usepackage{footnote}}
\makesavenoteenv{longtable}
\usepackage{graphicx}
\makeatletter
\newsavebox\pandoc@box
\newcommand*\pandocbounded[1]{% scales image to fit in text height/width
  \sbox\pandoc@box{#1}%
  \Gscale@div\@tempa{\textheight}{\dimexpr\ht\pandoc@box+\dp\pandoc@box\relax}%
  \Gscale@div\@tempb{\linewidth}{\wd\pandoc@box}%
  \ifdim\@tempb\p@<\@tempa\p@\let\@tempa\@tempb\fi% select the smaller of both
  \ifdim\@tempa\p@<\p@\scalebox{\@tempa}{\usebox\pandoc@box}%
  \else\usebox{\pandoc@box}%
  \fi%
}
% Set default figure placement to htbp
\def\fps@figure{htbp}
\makeatother





\setlength{\emergencystretch}{3em} % prevent overfull lines

\providecommand{\tightlist}{%
  \setlength{\itemsep}{0pt}\setlength{\parskip}{0pt}}



 


\KOMAoption{captions}{tableheading}
\makeatletter
\@ifpackageloaded{caption}{}{\usepackage{caption}}
\AtBeginDocument{%
\ifdefined\contentsname
  \renewcommand*\contentsname{Table of contents}
\else
  \newcommand\contentsname{Table of contents}
\fi
\ifdefined\listfigurename
  \renewcommand*\listfigurename{List of Figures}
\else
  \newcommand\listfigurename{List of Figures}
\fi
\ifdefined\listtablename
  \renewcommand*\listtablename{List of Tables}
\else
  \newcommand\listtablename{List of Tables}
\fi
\ifdefined\figurename
  \renewcommand*\figurename{Figure}
\else
  \newcommand\figurename{Figure}
\fi
\ifdefined\tablename
  \renewcommand*\tablename{Table}
\else
  \newcommand\tablename{Table}
\fi
}
\@ifpackageloaded{float}{}{\usepackage{float}}
\floatstyle{ruled}
\@ifundefined{c@chapter}{\newfloat{codelisting}{h}{lop}}{\newfloat{codelisting}{h}{lop}[chapter]}
\floatname{codelisting}{Listing}
\newcommand*\listoflistings{\listof{codelisting}{List of Listings}}
\makeatother
\makeatletter
\makeatother
\makeatletter
\@ifpackageloaded{caption}{}{\usepackage{caption}}
\@ifpackageloaded{subcaption}{}{\usepackage{subcaption}}
\makeatother
\usepackage{bookmark}
\IfFileExists{xurl.sty}{\usepackage{xurl}}{} % add URL line breaks if available
\urlstyle{same}
\hypersetup{
  pdftitle={Divergent Post-Out-of-Africa Metabolic Adaptations and a Miscalibrated Clinical Reference: The Evolutionary Basis of Global Type 2 Diabetes Disparities},
  pdfauthor={Emilio García-Morán},
  pdfkeywords={type 2 diabetes, ancestral metabolic
architecture, LDHA, LDHC, eQTL, paleogenomics, Beringia, Neolithic
transition, precision medicine, admixture},
  colorlinks=true,
  linkcolor={blue},
  filecolor={Maroon},
  citecolor={Blue},
  urlcolor={Blue},
  pdfcreator={LaTeX via pandoc}}


\title{Divergent Post-Out-of-Africa Metabolic Adaptations and a
Miscalibrated Clinical Reference: The Evolutionary Basis of Global Type
2 Diabetes Disparities}
\author{Emilio García-Morán}
\date{17 May 2026}
\begin{document}
\maketitle


\thispagestyle{empty}

\vspace{1cm}

\begin{center}
\textcolor[HTML]{2E6DA4}{\textit{Target journal: Genome Medicine}}\\[0.3cm]
\texttt{emilio.garcia.moran@uva.es} \quad · \quad ORCID: 0000-0002-2487-6686\\[0.2cm]
\textit{Data \& code:} \url{https://github.com/EmilioGarciaMoran/EastWestDM}\\
\textit{Zenodo DOI:} 10.5281/zenodo.20260486\\[0.2cm]
\textcolor[HTML]{666666}{\textit{Competing interests: None declared · Funding: None declared}}
\end{center}

\hrule
\vspace{0.5cm}

\subsection*{Abstract}\label{abstract}
\addcontentsline{toc}{subsection}{Abstract}

Type 2 diabetes (T2D) disproportionately affects East Asian (EAS),
admixed American (AMR), and African-ancestry populations compared to
non-Finnish Europeans (NFE). We propose that this disparity reflects not
differential pathology but the application of a clinical reference
framework calibrated on one post-Neolithic European metabolic
architecture to populations whose equally successful evolutionary
trajectories diverged prior to --- and were never subject to --- the
selective pressures that drove European metabolic remodelling.

We analysed 168 Eastern and 67 Western variants from four archaic and
early modern human genomes --- Denisovan D17, Ust'-Ishim
(\textasciitilde45 ka), Vindija Neandertal (\textasciitilde44 ka), and
Ötzi (\textasciitilde5.3 ka) --- against gnomAD v4. The ancestral NCDN
variant (chr1:35,565,741) is present at \textasciitilde70\% in
sub-Saharan African populations (San 92\%, Biaka 91\%), confirming
African ancestral origin predating the Out-of-Africa dispersal. Vindija
Neandertal already carries Western regulatory variants prior to
agriculture (\textasciitilde44 ka), indicating that the European
divergence began in pre-Neolithic temperate environments. Unadmixed
Native American populations (Surui, Karitiana: 100\%; Pima: 77\%; Maya:
67\%) and EAS (78\%) carry the ancestral architecture at or above the
African baseline, reflecting cold-environment enrichment during the
Beringian migration (\textasciitilde15,000 ya). NFE Europeans carry the
variant at 3\% --- a 23-fold reduction from the African baseline
consistent with Neolithic counter-selection. A single upstream variant
at chr11:18,411,103 (C\textgreater T) acts as a coordinated trans-eQTL
suppressing LDHC expression across seven tissues in GTEx v8 (skeletal
muscle \emph{p} = 3.71 × 10\textsuperscript{−64}, NES = −0.935;
independent replication in eQTLGen \emph{n} = 31,684, Z = 75.96,
\emph{p} = 3.27 × 10\textsuperscript{−310}). SLC2A3 (GLUT3) remains
invariant across 430,000 years of hominin evolution, establishing brain
glucose transport as a non-negotiable evolutionary constraint. African
Americans carry intermediate frequencies (ASW 49\%, ACB 73\%) consistent
with post-colonial admixture superimposed on the African ancestral
baseline. These findings reframe global T2D disparities as a consequence
of a miscalibrated clinical reference standard, with immediate
implications for ancestry-informed diagnostic thresholds and
pharmacogenomic stratification.

\textbf{Keywords:} type 2 diabetes · ancestral metabolic architecture ·
LDHA · LDHC · eQTL · paleogenomics · Beringia · Neolithic transition ·
precision medicine · admixture

\hrule
\vspace{0.3cm}

\subsection{Introduction}\label{introduction}

The global burden of type 2 diabetes is distributed with striking
population specificity. East Asian and admixed American populations show
higher prevalence, earlier onset, greater severity at equivalent body
mass index, and differential response to dietary and pharmacological
interventions compared to Europeans. African-ancestry populations share
this elevated burden and are additionally underrepresented in
genome-wide association studies and pharmacogenomic trials. The three
population groups most poorly served by current clinical frameworks ---
EAS, AMR, and African-ancestry --- thus share a common feature:
systematic underperformance of risk models calibrated on European
reference cohorts.

The thrifty genotype hypothesis posited a single ancestral adaptation to
feast-famine cycles. The thrifty phenotype hypothesis shifted causality
to developmental programming. Neither framework explains the shared
elevated burden across three geographically and historically distinct
population groups, nor does either account for why non-Finnish Europeans
--- who also experienced prehistoric food scarcity --- specifically do
not share it.

We propose a framework centred on a universal ancestral metabolic
architecture --- present in Africa at \textasciitilde70\% frequency
prior to the Out-of-Africa dispersal --- that subsequently diverged
along two distinct post-Out-of-Africa trajectories under radically
different selective pressures. The clinical consequences are a product
of the encounter between these trajectories and a medical reference
standard calibrated, by historical circumstance, on only one of them.

The metabolic architectures documented here represent independent
adaptive solutions of equivalent fitness. The European lineage,
beginning with pre-agricultural Neandertals in temperate environments
and consolidated by the Neolithic transition, optimised for carbohydrate
processing from plant-based substrates. The Siberian and Beringian
lineage optimised for maximal glycolytic efficiency under cold stress
and caloric scarcity, with brain glucose supply as the inviolable
constraint in both cases. The clinical problem arises not from the
architectures themselves but from their encounter with a globally
uniform food environment and a medical reference standard that, by
historical accident, reflects only one of them.

\subsection{Results}\label{results}

\subsubsection{Eastern and Western variant
discovery}\label{eastern-and-western-variant-discovery}

Comparison of D17 Denisovan and Ust'-Ishim (\textasciitilde45 ka,
Siberia) genomes against gnomAD v4 identified 168 high-confidence
Eastern variants after C→T/G→A damage filtering (mapDamage 2.2.1;
PMDtools threshold ≥3). Comparison of Vindija Neandertal and Ötzi
genomes identified 67 Western variants. Pathway analysis revealed
Eastern enrichment in glycolysis regulators (PFKM, LDHA, LDHC), adipose
biology (ADIPOQ, FASN), and cold thermosensation (TRPA1); Western
enrichment in lipid metabolism (ACACA), thermogenesis (PPAR$\gamma$GC1A), and
plant-based nutrient processing. Eighty-six percent of Eastern variants
showed symmetric polarisation toward AMR/EAS populations (Fisher
\emph{p} = 3.2 × 10\textsuperscript{−18}), with LDHA/LDHC regulatory
regions showing 21.4-fold enrichment.

\subsubsection{SLC2A3 invariance across 430,000
years}\label{slc2a3-invariance-across-430000-years}

Pileup analysis of five Sima de los Huesos specimens (PRJEB10597;
remapped to hg38; mean depth 1.18×) at SLC2A3 exonic and promoter
positions reveals no fixed differences relative to modern humans at 11
canonical positions examined. Combined with conservation data from
Vindija, D17, and Ust'-Ishim, SLC2A3 has remained functionally invariant
for at least 430,000 years across hominin lineages, establishing brain
glucose transport as a non-negotiable evolutionary constraint around
which both metabolic architectures were independently organised.

\subsubsection{LDHA/LDHC trans-eQTL across seven
tissues}\label{ldhaldhc-trans-eqtl-across-seven-tissues}

GTEx v8 analysis identified chr11:18,411,103 (C\textgreater T;
rs7951953) --- an Eastern-specific upstream variant 1,228 bp from the
LDHC transcription start site --- as a coordinated trans-eQTL
suppressing LDHC expression in seven tissues simultaneously (Table 1).
Effect sizes ranged from NES = −0.935 (skeletal muscle) to NES = −0.620
(whole blood), all in the same negative direction. The Eastern-specific
allele (C\textgreater T) is associated with lower LDHC expression in all
seven tissues, consistent with a regulatory variant that suppresses LDHC
transcription, shifting metabolic flux toward oxidative phosphorylation
--- thermogenically advantageous under cold stress, potentially
impairing metabolic flexibility under caloric excess. Independent
replication in eQTLGen (\emph{n} = 31,684; Z = 75.96; \emph{p} = 3.27 ×
10\textsuperscript{−310}) confirms this as among the most statistically
robust regulatory signals in the human genome.

\begin{longtable}[]{@{}llll@{}}
\caption{LDHA/LDHC trans-eQTL across GTEx v8 tissues (chr11:18,411,103
C\textgreater T). NES: normalised effect size (negative = lower LDHC
expression in Eastern allele carriers). eQTLGen: independent whole-blood
replication cohort.}\label{tbl-eqtl}\tabularnewline
\toprule\noalign{}
Tissue & \emph{p}-value & NES & \emph{n} \\
\midrule\noalign{}
\endfirsthead
\toprule\noalign{}
Tissue & \emph{p}-value & NES & \emph{n} \\
\midrule\noalign{}
\endhead
\bottomrule\noalign{}
\endlastfoot
Skeletal muscle & 3.71 × 10\textsuperscript{−64} & −0.935 & 706 \\
Subcutaneous adipose & 1.46 × 10\textsuperscript{−59} & −0.881 & 581 \\
Visceral adipose & 1.11 × 10\textsuperscript{−41} & −0.812 & 469 \\
Whole blood & 7.39 × 10\textsuperscript{−30} & −0.620 & 670 \\
Pancreas & 7.33 × 10\textsuperscript{−26} & −0.751 & 305 \\
Liver & 9.12 × 10\textsuperscript{−24} & −0.734 & 208 \\
Brain cortex & 2.70 × 10\textsuperscript{−22} & −0.698 & 205 \\
eQTLGen (blood) & 3.27 × 10\textsuperscript{−310} & Z = 75.96 &
31,684 \\
\end{longtable}

The brain cortex eQTL (\emph{p} = 2.70 × 10\textsuperscript{−22})
alongside SLC2A3 invariance establishes a critical dissociation: brain
glucose \emph{transport} was conserved for 430,000 years, while brain
glucose \emph{metabolism} retains population-specific regulatory
variation.

\subsubsection{Global frequency distribution: an ancestral variant with
asymmetric selective
history}\label{global-frequency-distribution-an-ancestral-variant-with-asymmetric-selective-history}

Query of the gnomAD v3.1.2 HGDP+1000G joint callset at the NCDN locus
(chr1:35,565,741) revealed a global frequency distribution consistent
with an ancestral allele present prior to the Out-of-Africa dispersal,
subsequently subject to asymmetric selection across lineages (Table 2,
Figures 5--6).

Sub-Saharan African populations carry the variant at a mean frequency of
70.4\% (range 49--92\%), establishing the ancestral baseline. San and
Biaka Pygmy populations show the highest frequencies (91.7\% and 90.9\%
respectively), consistent with long-term maintenance of the ancestral
state. EAS (78\%) and unadmixed Native American populations
(Surui/Karitiana: 100\%; Pima: 77\%; Maya: 67\%) show frequencies at or
above the African baseline. NFE Europeans show near-complete loss of the
variant (3\%). African Americans carry intermediate frequencies (ASW:
49\%; ACB: 73\%) consistent with variable European admixture
superimposed on the African ancestral baseline. The independent
replication SNP at chr1:35,542,512 shows consistent directional effect
across all populations.

\begin{longtable}[]{@{}llll@{}}
\caption{Global frequency of ancestral NCDN variant (chr1:35,565,741) by
population. Sources: HGDP, 1000 Genomes Project, gnomAD
v3.1.2.}\label{tbl-freq}\tabularnewline
\toprule\noalign{}
Population & \emph{n} & Frequency & Selective context \\
\midrule\noalign{}
\endfirsthead
\toprule\noalign{}
Population & \emph{n} & Frequency & Selective context \\
\midrule\noalign{}
\endhead
\bottomrule\noalign{}
\endlastfoot
San (Kalahari) & 12 & 92\% & Ancestral African baseline \\
Biaka Pygmy (Congo) & 44 & 91\% & Ancestral African baseline \\
Mbuti Pygmy (Congo) & 24 & 83\% & Ancestral African baseline \\
Surui (Amazonia) & 14 & 100\% & Cold-environment enrichment \\
Karitiana (Amazonia) & 20 & 100\% & Cold-environment enrichment \\
EAS (gnomAD) & 5,174 & 78\% & Cold-environment enrichment \\
Pima (N. Mexico) & 22 & 77\% & Cold-environment enrichment \\
Yoruba (Nigeria) & 40 & 70\% & Ancestral African baseline \\
ACB (African Caribbean) & 188 & 73\% & Admixture dilution \\
Maya (Mesoamerica) & 36 & 67\% & Cold-environment enrichment \\
AFR gnomAD (global) & 41,322 & 56\% & Pan-African mean \\
Peruvian (Andes) & 170 & 53\% & Low European admixture \\
ASW (African American) & 104 & 49\% & Admixture dilution \\
Mexican (MXL, 1000G) & 126 & 39\% & Post-1521 European admixture \\
AMR gnomAD & 15,244 & 31\% & Post-1521 European admixture \\
Colombian (CLM, 1000G) & 188 & 29\% & Post-1521 European admixture \\
NFE (gnomAD) & 67,976 & 3\% & Neolithic counter-selection \\
\end{longtable}

\subsection{Discussion}\label{discussion}

\subsubsection{Two adaptive solutions of equivalent
fitness}\label{two-adaptive-solutions-of-equivalent-fitness}

The metabolic architectures documented here represent independent
evolutionary solutions of equivalent adaptive success --- not
differential pathology. The European lineage, beginning with
pre-agricultural Neandertals in temperate environments and consolidated
by the Neolithic transition, optimised for carbohydrate processing from
plant-based substrates: ACACA lipogenesis, PPAR$\gamma$GC1A thermogenesis, PPAR$\gamma$γ
adipogenesis calibrated for cereal-derived glucose. The Siberian and
Beringian lineage optimised for maximal glycolytic efficiency under cold
stress and caloric scarcity, with brain glucose supply as the inviolable
constraint. The clinical problem arises not from the architectures
themselves but from their encounter with a globally uniform food
environment and a medical reference standard that, by historical
accident, reflects only one of them.

\subsubsection{The European divergence began before
agriculture}\label{the-european-divergence-began-before-agriculture}

Vindija Neandertal (\textasciitilde44 ka) already carries Western
regulatory variants in the absence of cereals, domesticated animals, or
settled lifestyles --- predating the Neolithic transition by
approximately 35,000 years. The temperate European environment,
characterised by moderate cold stress and reliance on megafaunal fat as
caloric substrate, appears to have initiated counter-selection against
the ancestral African glycolytic architecture prior to the Neolithic.
The Neolithic transition accelerated a process already underway. Ötzi
(5,300 ya) represents not the origin of the Western metabolic
architecture but its completion.

Whether the ancestral NCDN architecture was present in Sima de los
Huesos hominin lineages (430 ka) and subsequently lost during the
European divergence remains undetermined. Coverage at chr1:35,565,741 in
SH specimens was insufficient for allelic inference (0 reads; mean
genome-wide depth 1.18×). This question represents a compelling target
for higher-coverage ancient European genome sequencing.

\subsubsection{From regulatory variant to glycolytic
mechanism}\label{from-regulatory-variant-to-glycolytic-mechanism}

The identification of chr11:18,411,103 (C\textgreater T) as a
coordinated trans-eQTL suppressing LDHC expression across seven
metabolically relevant tissues converts this population genomic
observation into a tractable molecular mechanism. LDHC encodes the
muscle/testis isoform of lactate dehydrogenase, catalysing the final
step of anaerobic glycolysis. Its coordinated downregulation across
muscle, adipose, pancreas, liver, and brain cortex --- all at negative
NES --- is consistent with a pleiotropic regulatory programme rather
than a tissue-specific effect. In calorically abundant environments,
LDHC suppression shifts the pyruvate/lactate equilibrium toward
oxidative phosphorylation --- thermogenically advantageous under Arctic
conditions, potentially impairing metabolic flexibility under caloric
excess.

\subsubsection{Regulatory ghost
adaptation}\label{regulatory-ghost-adaptation}

LDHA missense variants shared between D17 and Ust'-Ishim are absent from
gnomAD v4, indicating active purifying selection on coding changes. Yet
the surrounding regulatory landscape persists at high frequency across
most human populations outside Europe. We term this pattern
\emph{regulatory ghost adaptation}: coding changes eliminated by
purifying selection, regulatory architecture co-selected with them
persisting in non-coding space, visible only through eQTL analysis. The
European loss of this architecture --- rather than its presence
elsewhere --- is the evolutionary anomaly requiring explanation.

\subsubsection{The Beringia bottleneck: cold selection as amplifier, not
origin}\label{the-beringia-bottleneck-cold-selection-as-amplifier-not-origin}

The peopling of the Americas did not introduce a novel metabolic
architecture --- it exported and amplified an ancestral African one.
Unadmixed Native American populations carry the Eastern NCDN variant at
67--100\%, at or above the African baseline, consistent with positive
selection during the Arctic crossing (\textasciitilde15,000 ya). The
stepwise frequency reduction through populations of increasing European
admixture --- Peruvians (53\%), Mexicans (39\%), AMR gnomAD (31\%),
Colombians (29\%) --- constitutes a quantitative genomic record of
post-1521 demographic change. The same dilution mechanism operates in
African Americans (ASW: 49\%; ACB: 73\%), where post-colonial European
admixture has superimposed the European low-frequency architecture on
the African ancestral baseline. The \textasciitilde24\% frequency
difference between ASW and ACB reflects differential admixture
proportions, not differential selection.

\subsubsection{Three clinically undercharacterised populations, one
shared
mechanism}\label{three-clinically-undercharacterised-populations-one-shared-mechanism}

The populations most poorly served by European-calibrated clinical
thresholds share a single genomic feature: they carry the ancestral
metabolic architecture that Europeans specifically lost. SGLT2
inhibitors, GLP-1 receptor agonists, and thiazolidinediones were
validated predominantly in European-ancestry cohorts. The LDHC
trans-eQTL signal in pancreas (\emph{p} = 7.33 ×
10\textsuperscript{−26}), skeletal muscle (\emph{p} = 3.71 ×
10\textsuperscript{−64}), and liver (\emph{p} = 9.12 ×
10\textsuperscript{−24}) identifies specific tissues where differential
glycolytic flux may alter pharmacodynamic response. Stratification by
LDHA/LDHC regulatory haplotype --- rather than continental ancestry
label --- would recover this signal in admixed populations currently
misclassified by European-reference panels. The BMI thresholds
documented by Caleyachetty et al.~are not arbitrary measurement
failures: they are the epidemiological shadow of the evolutionary
divergence documented here.

\subsubsection{The human singularity and its metabolic
price}\label{the-human-singularity-and-its-metabolic-price}

SLC2A3 invariance across 430,000 years establishes brain glucose
priority as the non-negotiable constraint around which all hominin
metabolic architectures were organised. The periphery adapted ---
differently in different lineages, under different selective pressures,
across different timescales. The brain never compromised. The clinical
consequence is that the organ most dependent on stable glucose supply is
protected by the most conserved gene in our dataset, while the organs
that regulate glucose availability --- muscle, liver, pancreas, adipose
--- carry the full burden of evolutionary divergence.

\subsubsection{Limitations}\label{limitations}

This analysis is limited by available high-coverage archaic genomes. The
conservative C→T/G→A damage filter may remove genuine variants;
PMDtools-based filtering gave comparable directional results. Coverage
at the NCDN locus in Sima de los Huesos specimens was insufficient for
allelic inference (0 reads; mean depth 1.18×). Native American and
African HGDP sample sizes are small (\emph{n} = 12--44); replication in
the PAGE study, Arctic Biobank, and Mexico City Prospective Study is
required before clinical translation. The eQTLGen replication cohort
consists of whole blood; however, consistency of negative NES across all
seven GTEx tissues supports a pleiotropic regulatory programme. The
Neolithic counter-selection hypothesis requires formal testing with
time-series ancient DNA spanning the European Mesolithic-Neolithic
transition. CRISPR functional validation of the LDHA/LDHC regulatory
variant is needed to establish causality.

\subsection{Methods}\label{methods}

\subsubsection{Archaic genome processing and variant
calling}\label{archaic-genome-processing-and-variant-calling}

Five Sima de los Huesos specimens (PRJEB10597; Meyer et al.~2016) were
remapped to hg38 using AdapterRemoval (--qualitymax 93) → BWA-MEM (-k16)
→ Picard MarkDuplicates → mapDamage 2.2.1 (--rescale, MCMC). Variants
were called with GATK HaplotypeCaller. D17 Denisovan, Ust'-Ishim,
Vindija Neandertal, and Ötzi genomes were processed with identical
parameters. C→T and G→A variants were excluded at positions with
PMDtools score \textless3.

\subsubsection{Population frequency
analysis}\label{population-frequency-analysis}

gnomAD v4 allele frequencies were extracted via the gnomAD REST API.
Population-specific frequencies for HGDP and 1000 Genomes populations
were extracted via remote bcftools query from the gnomAD v3.1.2 HGDP+TGP
joint callset. Symmetric polarisation was assessed using Fisher's exact
test with Bonferroni correction. African population frequencies were
extracted for 15 sub-Saharan populations including deep-branching
lineages (San, Biaka Pygmy, Mbuti Pygmy).

\subsubsection{eQTL analysis}\label{eqtl-analysis}

GTEx v8 single-tissue eQTL data were queried for all Eastern variants
within 1 Mb of LDHC (ENSG00000166796). The eQTLGen cis-eQTL dataset
(2019-12-11 release; \emph{n} = 31,684; FDR \textless{} 0.05) was
queried for independent replication. Effect size direction (NES) was
compared across all tissues to assess pleiotropic consistency.

\subsubsection{SLC2A3 conservation
analysis}\label{slc2a3-conservation-analysis}

SLC2A3 exonic positions were queried in all archaic genomes using
samtools mpileup (minimum base quality 20, minimum mapping quality 25).
Positions covered at ≥1× depth were considered for conservation
analysis.

\subsubsection{Statistical analysis and figure
generation}\label{statistical-analysis-and-figure-generation}

All analyses were performed in R v4.5.2 (ggplot2, dplyr) and Python 3.11
(matplotlib, geopandas, pandas). All code is available at
https://github.com/EmilioGarciaMoran/EastWestDM.

\subsection{Data Availability}\label{data-availability}

All analysis code, variant lists, population frequency tables, and
figure generation scripts are available at:
\url{https://github.com/EmilioGarciaMoran/EastWestDM} (DOI:
10.5281/zenodo.20260486). Raw archaic genome data: PRJEB10597 (Sima de
los Huesos). gnomAD: \url{https://gnomad.broadinstitute.org}. GTEx:
\url{https://gtexportal.org}. eQTLGen: \url{https://eqtlgen.org}.

\subsection{Figure Legends}\label{figure-legends}

\textbf{Figure 1.} East-West variant frequency polarisation. 86\% of
Eastern variants enriched in AMR/EAS vs NFE (Fisher \emph{p} = 3.2 ×
10\textsuperscript{−18}). Bars: proportion of Eastern (blue) and Western
(red) variants in each population group.

\textbf{Figure 2.} Pathway enrichment. Eastern: glycolysis (PFKM, LDHA,
LDHC), adipose (ADIPOQ, FASN), cold thermosensation (TRPA1). Western:
lipogenesis (ACACA), thermogenesis (PPAR$\gamma$GC1A).

\textbf{Figure 3.} SLC2A3 invariance across 430,000 years. Pileup at 11
canonical positions across five Sima de los Huesos specimens, D17,
Ust'-Ishim, Vindija, Ötzi. No fixed differences observed.

\textbf{Figure 4.} LDHA/LDHC trans-eQTL. Panel A: NES by tissue (GTEx
v8). Panel B: eQTLGen replication (\emph{n} = 31,684; Z = 75.96;
\emph{p} = 3.27 × 10\textsuperscript{−310}). Panel C: regional
association plot chr11:18.35--18.55 Mb.

\textbf{Figure 5.} Global frequency gradient of ancestral NCDN variant
(chr1:35,565,741). Bars ordered by frequency. Colours: deep-branching
African (dark red), unadmixed Native American/EAS (red), admixed African
American/AMR (amber), European (blue). Open circles: replication at
chr1:35,542,512.

\textbf{Figure 6.} World map of ancestral NCDN variant frequency. Circle
size proportional to allele frequency; colour scale: blue (NFE 3\%) to
red (Surui/Karitiana 100\%). Arrows: Out-of-Africa (\textasciitilde70
ka, black); Beringia cold-environment enrichment (\textasciitilde15 ka,
dark red); Neolithic counter-selection in European lineages
(\textasciitilde8--5 ka, blue); post-colonial admixture events (dashed
grey). Source: HGDP + 1000 Genomes + gnomAD v3.1.2.

\begin{center}\rule{0.5\linewidth}{0.5pt}\end{center}

\emph{``The periphery adapted; the brain never compromised.''}

\begin{center}\rule{0.5\linewidth}{0.5pt}\end{center}

\vspace{0.5cm}

\textbackslash begin\{center\}
\textcolor[HTML]{666666}{\small EastWestDM v7.1 · García-Morán 2026 · Universidad de Valladolid · DOI: 10.5281/zenodo.20260486}
\textbackslash end\{center\textgreater{}




\end{document}
